\Needspace{18\baselineskip}
\Article{19}
Damit stehen zwei Methoden zur Verfügung, die sich für die praktische Durchführung besonders gut eignen. Die folgenden Tabellen zeigen die Anordnungen im Einzelnen.

\expandafter\def\csname HistoricalNoteText11\endcsname{Genauer gesagt, senkrecht zur vertikalen Ebene, der der Wert $u=u^0$ entspricht, also zu der Ebene, in der sich die magnetische Achse bei fehlender zweiter Nadel im Gleichgewicht befindet. In der Praxis kann man den Unterschied jedoch stets unbedenklich vernachlässigen, sowohl weil er klein ist als auch aufgrund eben der Überlegung, von der wir im vorangehenden Artikel unter III ausgegangen sind.}
\begin{center}
\textit{Erste Methode.}

Mittelpunkt und Achse der zweiten Nadel liegen beide auf einer Geraden senkrecht zum magnetischen Meridian\HistoricalNote{11}.

\begin{tabular}{c|c|c|l|l}
& & &Mitte&Nordpol\\
Ablenkung&\multicolumn{2}{c|}{Nadellage}&nach&nach\\\hline
$u=u'$&$\psi=90^0$&$U=90^0$&Osten&Osten\\
$u=u''$&$\psi=90^0$&$U=270^0$&Osten&Westen\\
$u=u'''$&$\psi=270^0$&$U=90^0$&Westen&Osten\\
$u=u''''$&$\psi=270^0$&$U=270^0$&Westen&Westen
\end{tabular}
\end{center}
\SourceLeaf{112}
\expandafter\def\csname HistoricalNoteText12\endcsname{Die bei diesen Versuchen verwendeten Nadeln sind etwa $0^{\mathrm m}3$ lang. Hätten wir es gewagt, den Term $R^{-7}$ in die Rechnungen einzubeziehen, wäre die Sicherheit eher gesunken als gestiegen.}
\Needspace{13\baselineskip}
\begin{center}
\textit{Zweite Methode.}\par
Der Mittelpunkt der zweiten Nadel liegt im magnetischen Meridian, ihre Achse steht senkrecht auf ihm.
\begin{tabular}{c|cc|cc}
Ablenkung&\multicolumn{4}{c}{Nadellage}\\
&&&Mitte nach&Nordpol nach\\\hline
$u=u'$&$\psi=0$&$U=270^0$&Norden&Westen\\
$u=u''$&$\psi=0$&$U=270^0$&Norden&Osten\\
$u=u'''$&$\psi=180^0$&$U=90^0$&Süden&Westen\\
$u=u''''$&$\psi=180^0$&$U=90^0$&Süden&Osten
\end{tabular}
\end{center}

Setzen wir dann $\frac14(u'-u''+u'''-u'''')=v$ und
\[
\operatorname{tang}v=LR^{-(n+1)}+L'R^{-(n+3)}+L''R^{-(n+5)}+\text{usw.}
\]
so ergibt sich
\[
\begin{aligned}
\text{erste Methode}\qquad L&=\frac{nmM}{mT+\theta}\\
\text{zweite Methode}\qquad L&=\frac{mM}{mT+\theta}
\end{aligned}
\]
